\scanpage{4}
\mspage{2}
\deleted{\unsure{Deux} \unsure{1)} Un ensemble $L$ (l}

1) Deux ensembles de base $\unsure{L|L_*}$ (ensemble des \uline{lieux} du modèle) et $S$ (ensemble des \uline{segments} du modèle)

2) Une application $\unsure{S}\longrightarrow\mathcal P(L)$, $I-\tilde I$ (lieux sur un \unsure{segment}) -- i.e. une relation entre $S$ et $L$.

3) Une application $\delta:S\longrightarrow\mathcal P_2(L)$, \inserted{$I\longmapsto\delta I$,} telle que $\delta I\subset\tilde I$.

\deleted{\unsure{1}} $+$ axiomes

\uline{NB} J'ignore s'il faut supposer que $I$ est connu, quand on connaît $\tilde I$ et $\delta I$, i.e. s'il faut supposer $I\longrightarrow\mathcal P(L)\times\mathcal P_2(L)$ injectif.

\margin{7.6. On verra que $I$ est connu quand on connaît $\tilde I$, on identifie $I$ \unsure{à $\tilde I$}, et $S$ à une partie de $\mathcal P(L)$.}

Si $I$ est un segment, les éléments de $\delta I$ sont les lieux extrémités, et ceux de $\tilde I\setminus\delta I$ \inserted{$=\deleted{\unsure{\delta I}}\tilde I^\circ$} sont les lieux intérieurs au segment.

Un lieu est dit \uline{régulier} s'il est intérieur à \inserted{(au moins)} un segment, \uline{singulier} (\deleted{\unsure{ou}} critique, lieu de branchement) sinon.
\[
L_r=\bigcup_{I\in S}(\tilde I\setminus\delta I),\qquad L_s=L\setminus L_r
\]
Ainsi, pour tout $I\in S$, on a
\[
\tilde I\cap L_s\subset\delta I
\]
\mspage{3}

\rotatebox{90}{\reflectbox{\unsure{EXTRA STRONG}}}
\rotatebox{90}{\reflectbox{\unsure{EXTRA STRONG}}}
